\documentclass[letter]{res}

\setlength{\textheight}{9.5in}
\sectionskip=2.5ex plus 1ex minus .2ex
%\usepackage{fontspec}
\usepackage{geometry}
\usepackage{xcolor}
\usepackage[utf8]{inputenc}
\geometry{
 a4paper,
 total={210mm,297mm},
 left=20mm,
 right=30mm,
 top=12mm,
 bottom=12mm,
 bindingoffset=0mm
 }\usepackage{fontenc}
\usepackage{enumitem}
\setlist{nolistsep}
\usepackage[colorinlistoftodos]{todonotes}
\usepackage{hyperref}

%\setmainfont[Mapping=tex-text]{Linux Libertine O}


\begin{document}

\name{Jacob Merle Gershon\\[14pt]}
\address{ \href{https://jacobgershon.com}{jacobgershon.com} \\}
\address{ \texttt{nohsr3g@gmail.com}\\ }

\begin{resume}
\hspace{-8mm}\noindent\makebox[\linewidth]{\rule{\paperwidth}{0.4pt}}
% \noindent\makebox[\linewidth]{\rule{\paperwidth}{0.4pt}}
  % \makebox[\linewidth]{\rule{\paperwidth}{0.4pt}}


\section{Experience}
Prescient Design, Genentech \hfill Feb 2026 - Present\\
{\sl Machine Learning Scientist}\\
 \vspace{-2mm}
 \begin{itemize}
 \item Post-training structure prediction models.
 \item Building agentic workflows for drug discovery.
 \end{itemize}
  \vspace{-2mm}

\section{Education}
\textit{Ph.D.} - Molecular Engineering \hfill 2025\\
Institute for Protein Design, University of Washington\\
Dissertation: \textit{Protein Design at Library Scale}\\[3pt]
\textit{M.S.} - Biomedical Engineering \hfill 2021\\
State University of New York at Binghamton\\
Thesis: \textit{Enhancing Formation of Stem-cell Derived 3D Cardiac Tissue with Extracellular Matrix}\\[3pt]
\textit{B.S.} - Biomedical Engineering \hfill 2020\\
State University of New York at Binghamton\\
Honors: \textit{Summa Cum Laude}
  \vspace{-2mm}

\section{Research}
Institute for Protein Design, David Baker \hfill 2021 - 2025\\
{\sl Graduate Research Assistant}\\
 \vspace{-2mm}
 \begin{itemize}
 \item  Developed a diffusion model in protein sequence space using a structure prediction framework.
 \item Incorporated experimental data to guide the model, demonstrating the power of an informed prior.
 \item Built a discrete Bayesian optimization workflow to enhance protein fitness via lab-in-a-loop.
 \item Designed scalable pipelines for creating custom-filtered libraries for high throughput screens.
 \item Trained autoencoders with novel variable rate compression strategies for protein-ligand complexes.
 \end{itemize}



Regenerative Therapeutics at Binghamton, Tracy Hookway \hfill 2018 - 2021\\
{\sl Undergraduate Research Assistant}\\
 \vspace{-2mm}
 \begin{itemize}
 \item Engineered 3D cardiac microtissues to enhance biophysical organization.
 \item Designed a bioreactor for growth of cardiac mesh tissues.
 \end{itemize}
   \vspace{-2mm}


\section{Publications \& Presentations}
Computational design of cysteine proteases \hfill 2025\\
\textit{Choi, H., Coventry, B., et al., Gershon, J., et al.,} \href{https://www.biorxiv.org/content/10.1101/2025.11.21.689808}{biorxiv.org/content/10.1101/2025.11.21.689808}\\[3pt]
Hierarchical Protein Backbone Generation with Latent and Structure Diffusion \hfill 2025\\
\textit{Yim, J.*, Jaakik, M.*, Liu, G.*, Gershon J.*, et al.,} \href{https://arxiv.org/abs/2504.09374}{https://arxiv.org/abs/2504.09374}\\[3pt]
Multistate and functional protein design using RoseTTAFold sequence space diffusion \hfill 2024\\
\textit{Lisanza, S.*, Gershon, J.*, et al.,} \href {https://www.nature.com/articles/s41587-024-02395-w}{nature.com/articles/s41587-024-02395-w}\\[3pt]
Towards Guiding Generative Models with Experimental Data \hfill 2024\\
\textit{ML Protein Engineering Seminar,} \href{https://youtu.be/hbkDL_vxGKE}{youtu.be/hbkDL\_vxGKE}\\[3pt]
% Exploring Sequence Space Diffusion Models \hfill 2023\\
% \textit{Rosettacon, Oral Presentation}\\
% \\
Enhancing Formation and Alignment of 3D Cardiac Microtissues \hfill 2021\\
\textit{Masters thesis,} State University of New York at Binghamton\\[3pt]
Construction of 3D Cardiac Microtissues with Extracellular Matrix Additives \hfill 2018\\
\textit{ Biomedical Engineering Society Annual Meeting, oral presentation}
    \vspace{-2mm}

\section{Service, Mentorship \& Teaching}
GEM Workshop on Generative and Experimental Biomolecular Design \hfill 2024 - Present\\
\textit{Co-organizer} - Workshops at ICLR and NeurIPS, \href{https://www.gembio.ai}{gembio.ai}\\[3pt]
Undergraduate Mentorship Program at University of Washington \hfill 2023 - 2025 \\
\textit{Mentor} - Trained students to process data and train sequence to function prediction models. \\[3pt]
Introduction to Structural Biology \hfill 2022\\
\textit{Teaching Assistant} - Led discussion sections and lectured unit on methods development\\
 \vspace{-4mm}

% Introduction to Biomaterials\\
% \textit{Teaching Assistant} - Helped with breakout sections during lectures and grading \hfill 2019\\
% \\
% Introduction to Biochemistry\\
% \textit{Teaching Assistant} - Worked with other TAs to organize recitation sections \hfill 2018\\
\section{Software}
    Python, PyTorch, Hydra, Biotite, NumPy, GitHub, Matplotlib, Seaborn
 \vspace{-2mm}

\end{resume}
\end{document}
